\chapter{Anthropic Assumptions}\label{app:anthropics}

Chapter 8 identified missing structure. $M$ is a semimeasure over program-output prefixes, not a probability distribution over observer occurrences. To ask which matching occurrence is \emph{you}, one must first define a reference class or sample space and then assign self-location weights. Cosmological claims about ordinary observers and fluctuations require still more: a physical model that produces the occurrences. This appendix sketches several additions and explains why the book adopts none.

The \emph{Self-Sampling Assumption} (SSA) says: reason as if your current observer-moment were sampled from a specified reference class, for instance observers with similar information or members of a species. Once the class and its weighting are fixed, SSA can update beliefs about your place within it. This is one route to the doomsday argument. Its answer depends on who enters the class and what ``sampled'' means.

The \emph{Self-Indication Assumption} (SIA) gives hypotheses more weight when they contain more eligible observers like you. It thereby couples existence evidence to population size, once eligibility and counting have been defined.

The assumptions can disagree on Sleeping Beauty, the doomsday argument, and fluctuation observers. SSA is sensitive to the reference class; SIA faces the ``presumptuous philosopher'' problem. A more general approach specifies a \emph{xerographic distribution}, an explicit first-person probability over otherwise matching occurrences. That notation makes the missing choice visible rather than pretending $M$ already made it. None follows from Bayes, Kraft, $M$, or dominance.

The \emph{speed prior}, due to Schmidhuber, changes the program prior by penalizing computation time as well as description length. It can downweight outputs that require long computations, at the cost of disfavoring computational depth. But it still weights programs or outputs. It does not by itself identify observer occurrences, define a reference class, or provide a complete solution to the Boltzmann-observer problem.

The proposals do not all do the same job. SSA, SIA, and xerographic distributions add indexical structure. The speed prior changes program weights. A cosmological model determines which physical occurrences are present. Combining them can produce a sharper answer, but each ingredient must be defended separately.

This book adopts none, not as a verdict against them but as a matter of discipline. A reader who wants a sharper answer must specify an ontology, a cosmology where relevant, an occurrence definition, and a self-location distribution rather than hiding those choices inside the phrase ``universal prior.''
